% chapters/09/sub-cpmk.tex
\begin{subcpmkbox}
Bab ini menuntaskan \SUBCPMK{2.2} secara utuh. Ketiga indikatornya diperiksa,
dan buktinya berupa diagram keadaan, program NPC bermesin keadaan, serta
catatan pengujian setiap keadaan dan perpindahannya.

\vspace{0.5em}
\textbf{\SUBCPMK{2.2} Mengimplementasikan \istilah{Finite State Machine}
(\istilah{idle}, \istilah{patrol}, \istilah{chase}, \istilah{attack/flee})
untuk NPC} --- ranah C3, \CPMK{2}, \CPL{03}.

\vspace{0.5em}
{\setlength{\tabcolsep}{4pt}\footnotesize%
\begin{tabularx}{\textwidth}{@{}C{1.9cm} L{4.0cm} L{2.5cm} Y@{}}
\toprule
\headrow \textbf{Indikator} & \textbf{Rumusan} & \textbf{Dibahas pada} & \textbf{Bukti} \\
\midrule
\ind{a} & Menggambar diagram \istilah{state}--transisi NPC &
Bagian~\ref{sec:diagram-keadaan} &
Diagram keadaan--transisi buatan sendiri, dengan empat keadaan dan setiap
perpindahan yang diberi syarat yang dapat diperiksa. \\
\addlinespace
\ind{b} & Mengimplementasikan transisi berbasis jarak dan \istilah{cooldown} &
Bagian~\ref{sec:menerapkan-fsm} &
Kode mesin keadaan, beserta penjelasan setiap syarat perpindahan dan pewaktu
yang dipakai. \\
\addlinespace
\ind{c} & Menguji setiap \istilah{state} dan transisinya secara terpisah &
Bagian~\ref{sec:menguji-keadaan} &
Catatan perpindahan keadaan beserta nomor bingkainya, hasil skenario pengujian,
dan pemeriksaan keabsahan seluruh perpindahan. \\
\bottomrule
\end{tabularx}}
\end{subcpmkbox}
